\documentclass[11pt]{article}
\usepackage[margin=1.1in]{geometry}
\usepackage{amsmath, amssymb, amsthm, mathtools}
\usepackage{enumitem}
\usepackage[colorlinks=true, linkcolor=blue, citecolor=blue]{hyperref}

\newlist{conditions}{description}{1}
\setlist[conditions]{style=unboxed, leftmargin=0pt, labelsep=0.35em,
                     itemsep=5pt, parsep=2pt, font=\normalfont}

\newtheorem{theorem}{Theorem}
\newtheorem{lemma}{Lemma}
\newtheorem{proposition}{Proposition}
\newtheorem{corollary}{Corollary}
\newtheorem{assumption}{Assumption}
\theoremstyle{definition}
\newtheorem{definition}{Definition}
\newtheorem{remark}{Remark}
\newtheorem{example}{Example}

\newcommand{\R}{\mathbb{R}}
\newcommand{\E}{\mathbb{E}}
\newcommand{\cA}{\mathcal{A}}
\newcommand{\cB}{\mathcal{B}}
\newcommand{\cP}{\mathcal{P}}
\newcommand{\cN}{\mathcal{N}}
\newcommand{\cT}{\mathcal{T}}
\newcommand{\ip}[2]{\left\langle #1, #2 \right\rangle}
\newcommand{\KL}{\mathrm{KL}}
\newcommand{\Lip}{\operatorname{Lip}}
\DeclareMathOperator*{\argmax}{arg\,max}
\DeclareMathOperator{\supp}{supp}
\DeclareMathOperator{\conv}{conv}
\DeclareMathOperator{\diag}{diag}
\DeclareMathOperator{\tr}{tr}
\newcommand{\GC}[1]{\textsf{GC-#1}}

\title{Convergence of Magnetic Mirror Descent with Gaussian-Mixture Policies\\
on the Class $\mathfrak{G}$\\
\large A theorem, and a proof}
\author{\texttt{cont\_tinkering} / \texttt{theory2}}
\date{August 2026}

\begin{document}
\maketitle

\begin{abstract}
This note discharges the principal open item of \texttt{game\_class.tex} (its
\S5.6, item~1, and Decision~1 of \texttt{paper\_outline.tex}): the composition of
the three input-to-state-stable blocks into a global convergence theorem for the
class $\mathfrak{G}(M,\Sigma,\cN,\rho)$. We state a convergence theorem
(Theorem~\ref{thm:main}) and prove it. Closing the argument costs \emph{one}
substantive new hypothesis --- that the star-concavity modulus $\rho$ is
\emph{linear}, $\rho(r) \ge \alpha r$ with $\alpha > 0$ possibly tiny
(hypothesis~(L)) --- together with a quantified form of the qualitative small-gain
condition (D) already in the class, and the explicit constants that (S)
already promises but does not name. Hypothesis~(L) is not an extra restriction on
which games qualify in any case of interest: on a compact cell it is implied by
continuity of $\rho$ together with nondegeneracy of the branch maximum, and it is
what supplies the uniform branch curvature that the scale block needs anyway.

The technical content is three observations. First, the mean update's step is
proportional to the tracking error ($\|\nabla Q^{\mathrm{reg}}(a)\| \le L\|a -
\xi_j\|$, because the branch point is a critical point), so the drift each block
injects into its neighbours is proportional to that block's own error --- which is
what converts the qualitative ISS estimates into a genuine contraction rather
than convergence to an $O(\eta)$ ball. Second, and this is the resolution of
\S5.6 item~5, the Fisher preconditioner $\sigma^2$ appears in the contraction
rate \emph{and} in every drift term, so it cancels out of every cycle gain: the
small-gain condition is scale-free and step-size-free, and the ``scale collapse
throttles transport'' objection is quantitatively void. Third, and this is the
resolution of \S5.6 item~1, the tension between ``small gain wants
$\tau^{\mathrm{cat}}$ large'' and ``low bias wants $\tau^{\mathrm{cat}}$ small''
is only apparent: the magnet refresh is a Bregman proximal-point iteration whose
fixed point is the Nash of $\bar A$ for \emph{every} $\tau^{\mathrm{cat}} > 0$,
so $\tau^{\mathrm{cat}}$ enters the rate and not the limit, and may be taken as
large as the gain condition requires.

The conclusion is the one advertised: means converge to branch points, scales to
the floor, weights to the Nash of the smoothed atom matrix, and the induced
policy pair is $\varepsilon$-Nash with $\varepsilon = O((s^{\min}_0)^2 +
(s^{\min}_1)^2)$ --- with, as a by-product, a sharper constant than the one
previously stated, and a linear rate. Section~\ref{sec:remaining} lists what is
\emph{still} not proved, which is now a shorter and more specific list.
\end{abstract}

\tableofcontents

\section{Roadmap: the idea of every step}\label{sec:roadmap}

This section is the proof in words. Each entry says what is \emph{obvious} or
imported about a step, and then what the step actually has to add --- which is
almost always the same thing: the object being chased is moving, and the work is
in charging that motion to something that vanishes.

\subsection*{The three ideas}

\begin{enumerate}[leftmargin=1.4em]
\item \textbf{Every block's step is proportional to its own error.} A component
  mean moves fast only when it is far from where it is going, because the branch
  point is a critical point of the landscape it is climbing
  (Lemma~\ref{lem:dom}). The same is true of the categorical head. So the
  disturbance each block injects into its neighbours is proportional to that
  block's own error, and ``everything is disturbing everything'' becomes a system
  that can contract to zero rather than to a ball of radius $O(\eta)$.
\item \textbf{The preconditioner cancels.} The Fisher factor $\sigma^2$ slows the
  mean update down as the scales collapse --- and slows down, by at least as much,
  every disturbance the mean update is trying to reject. Around any cycle of the
  interconnection it appears once upstairs and once downstairs
  (Lemma~\ref{lem:cyclefree}). So the stability condition contains no $\sigma$ and
  no step size: it is a statement about the \emph{game}.
\item \textbf{The categorical magnet temperature is free.} It sets how stiff the
  weight block's target is against the means moving underneath it, and it does
  not move the limit (Lemma~\ref{lem:weight}(iii)). So it can be turned up until
  the loop closes, at no asymptotic cost.
\end{enumerate}

\subsection*{The blocks}

There are four errors --- the two players' mean tracking errors $\tilde r_0,
\tilde r_1$ and their two categorical errors $\delta_0, \delta_1$ --- plus the two
scale variables, which are handled separately because they converge in finite
time rather than asymptotically. The lemmas come in four groups.

\paragraph{Group 1: what a branch point is
(Lemmas~\ref{lem:dom}--\ref{lem:branchdrift}).}
\begin{description}[leftmargin=0pt, style=unboxed, font=\normalfont\itshape]
\item[Lemma~\ref{lem:dom}.] First, a point of logic: the branch point is
  \emph{not} defined as a critical point. It is data of hypothesis~(C) --- a point
  of the cell, posited to exist, toward which the gradient leans from everywhere
  in that cell --- and criticality is a \emph{consequence}, proved here along with
  uniqueness and maximality. (Setting it up this way is what makes (C) one
  checkable inequality rather than three separate claims, and it is why (C) can be
  verified by counting modes on a scale ladder.) What the lemma then adds, and
  the reason it exists: since the gradient vanishes at $\xi_j$ and is
  $L$-Lipschitz, it is at most $L$ times the distance to $\xi_j$ everywhere in the
  cell --- so the mean's step shrinks as it arrives. That one line is what later
  separates ``converges'' from ``converges to an $O(\eta)$ ball'', and it is used
  in every subsequent estimate.
\item[Lemma~\ref{lem:curv}.] Obvious: a nondegenerate maximum is curved. What it
  adds: hypothesis~(L) --- star-concavity with a \emph{linear} modulus --- is
  exactly the statement that this curvature is bounded away from zero uniformly
  in the opponent and the scale. So the uniform curvature the scale block needs
  does not have to be assumed separately; it comes free with (L).
\item[Lemma~\ref{lem:branchdrift}.] Obvious: change the smoothing width and the
  peak of the smoothed landscape moves. What it adds: it moves by $O(1)$ per unit
  of $s^2/2$, not per unit of $\log s$. Since the algorithm updates $\log s$, the
  branch drift caused by a scale step carries a factor $\sigma^2$. That factor is
  banked here and spent in Lemma~\ref{lem:scalefree}.
\end{description}

\paragraph{Group 2: the means (Lemmas~\ref{lem:contract}--\ref{lem:mean}).}
\begin{description}[leftmargin=0pt, style=unboxed, font=\normalfont\itshape]
\item[Lemma~\ref{lem:contract}.] Obvious: preconditioned gradient ascent on a
  star-concave function converges. What it adds: the rate, written in the metric
  the algorithm actually uses --- anisotropic $\Sigma_k$, the $w_k$ preconditioner
  artifact of \GC{Remark~2} --- and the observation that the second-order term of
  the expansion is proportional to $r^2$ rather than constant, so the contraction
  is geometric and not merely a decrease.
\item[Lemma~\ref{lem:inv}.] Obvious: if you always step toward the branch point
  you never leave the cell. What it adds: the explicit radius $R_i$ within which
  that is true, which is the budget hypothesis~(G$_\sigma$) then has to fit the
  entire transient inside.
\item[Lemma~\ref{lem:mean}.] This is Lemma~\ref{lem:contract} with the target
  moving. The branch point moves for four reasons --- the opponent's means, the
  opponent's weights, either player's scales, and the magnet centre --- and
  (D$^+$) is what lets each be charged separately. Output: one inequality reading
  ``error contracts by a fixed fraction, plus however far the target moved''.
\end{description}

\paragraph{Group 3: the scales (Lemmas~\ref{lem:scale}--\ref{lem:scalefree}).}
\begin{description}[leftmargin=0pt, style=unboxed, font=\normalfont\itshape]
\item[Lemma~\ref{lem:scale}.] Obvious: at a maximum the Laplacian is negative, so
  blurring costs value, so the scale update pushes the scale down --- this is
  \GC{Lemma~3}, and it is why no annealing schedule is needed. What it adds: the
  same holds throughout a ball of radius $\delta^\kappa$ around the branch point,
  uniformly in the opponent. That is what lets the scales be shown to fall while
  the means have only \emph{approximately} arrived, and it converts an asymptotic
  statement into a finite arrival time at the floor --- after which the scale
  contributes exactly zero disturbance, which is what Step 4 of the theorem needs.
\item[Lemma~\ref{lem:scalefree}.] The objection it answers is \GC{\S5.6},
  item~5: as the scales collapse the mean update's rate falls by four orders of
  magnitude, so how can transport possibly survive? Answer: the disturbance that
  the scale collapse injects into the mean block falls by \emph{six} --- $\sigma^2$
  from Lemma~\ref{lem:branchdrift} and $\sigma^2$ from the preconditioner in the
  scale update itself, against $\sigma^2$ in the contraction. The ratio, which is
  all that matters, improves as the scales fall.
\end{description}

\paragraph{Group 4: the weights (Lemma~\ref{lem:weight}).}
\begin{description}[leftmargin=0pt, style=unboxed, font=\normalfont\itshape]
\item[Lemma~\ref{lem:weight}.] Obvious, and imported: for a \emph{fixed} payoff
  matrix, tabular MMD converges to the regularised equilibrium, and iterating the
  magnet refresh drives that to Nash. This is the paper's spine (\GC{Lemma~1}(i))
  and needs no hypothesis on the game at all. What the lemma adds is everything
  that follows from the matrix \emph{not} being fixed: $\bar A(t)$ changes every
  step because the means and scales move under it, so the statement needed is a
  \emph{lag} bound --- the weights sit behind their moving target by an amount
  proportional to the speed of the means, which by Lemma~\ref{lem:dom} is itself
  proportional to the means' own tracking error. That is part~(v), and it is the
  whole content; parts (i), (ii), (iv) are its three constants (how fast the gap
  closes, how far the target moves per unit of matrix perturbation --- the
  $O(1/\tau^{\mathrm{cat}})$ stiffness --- and how fast $w$ moves per unit of its
  own error). Part~(iv) exists because the opponent's weights tilt \emph{your}
  landscape: this block is a feedback edge, not a read-out, which is why a
  small-gain argument is needed rather than a sequential one.
\end{description}

\subsection*{Composition, and the limit}

\begin{description}[leftmargin=0pt, style=unboxed, font=\normalfont\itshape]
\item[Lemma~\ref{lem:sg}.] Obvious: four contracting blocks, each disturbing the
  others, still converge if the disturbances are small enough. What it adds:
  ``small enough'' in a form that can be checked --- the spectral radius of the
  $4\times4$ matrix of (disturbance)/(own rate) ratios --- together with the
  invariant-ball version needed for the transient phase, while the scales are
  still falling and the disturbance has not yet reached zero.
\item[Lemma~\ref{lem:cyclefree}.] The bookkeeping that makes the previous lemma's
  hypothesis usable: around any cycle, each $\eta_i$ and each $\sigma$ appears
  once in a numerator and once in a denominator. So (G) can be evaluated offline
  from game data, before choosing any hyperparameter.
\item[Lemma~\ref{lem:fp}.] The four errors going to zero says the configuration
  stops moving; it does not say \emph{where}. This lemma names the limit: the
  unique configuration that is self-consistent --- every mean at the branch point
  induced by the opponent's policy, every weight vector the Nash of the matrix
  those means induce. Existence and uniqueness come from the same constants as
  (G), because the consistency map is a contraction with the same gains. Then (M)
  says the true equilibrium atoms are an $O(s^2)$-\emph{approximate} fixed point
  of that map, so the actual fixed point is $O(s^2)$ away from them.
\item[Lemma~\ref{lem:exp}.] Once the limit is located, exploitability is a short
  computation, and it comes out cleaner than expected: because the categorical
  head equalises all $M_0$ cells exactly at a Nash of $\bar A$, the regret is not
  a weighted average of per-atom smoothing deficits but literally
  $\max Q_\nu - \max (Q_\nu)_s$ --- a deficit at a single point.
\end{description}

\subsection*{The theorem}

Theorem~\ref{thm:main} is then six steps, in this order:
assign components to cells (K); assemble the four blocks into the comparison
system and apply Lemma~\ref{lem:sg}; note that the transient never leaves the
cells and enters the ball where the scales are monotone; let the scales fall to
the floor in finite time, at which point the exogenous disturbance is
\emph{exactly} zero; then the same small-gain estimate with zero input drives all
four errors to zero geometrically; finally identify the limit with
Lemma~\ref{lem:fp} and read off the bias with Lemma~\ref{lem:exp}.

The ordering of hypotheses matters and is not circular: $\tau^{\mathrm{cat}}$ is
chosen from the gain condition \eqref{eq:cycleweight}, then the refresh period $T$
from (T), then the step sizes from \eqref{eq:step}. This is the ordering
\GC{\S5.6}, item~1 could not settle, and Lemma~\ref{lem:weight}(iii) is what
unblocks it.

The corollaries are variations on the same argument:
Corollary~\ref{cor:unfloored} releases the floor and observes that the two
players' scales fall at proportional rates, so Lemma~\ref{lem:cyclefree} still
applies with time-varying rates whose sum diverges --- convergence survives,
degraded from geometric to polynomial. Corollary~\ref{cor:overcap} notes that
spare components are harmless because there is no spurious attractor inside the
class to catch them. Corollary~\ref{cor:flat} is the case $\kappa_i = 0$, where
Group~3 fails and nothing else does, so the conclusion degrades to convergence of
the induced measure. Proposition~\ref{prop:cap} evaluates (G) in closed form when
the coupling has finite rank.

\section{Setting, and what is imported}\label{sec:setting}

Throughout, references of the form \GC{Definition~1}, \GC{Lemma~1},
\GC{(2.4)} point into \texttt{theory2/game\_class.tex}; that document's notation
is used without change. We recall only what the proof consumes.

\paragraph{Game and equilibrium.} Player $0$ maximises, player $1$ minimises a
bounded measurable $u : \cA \times \cB \to \R$ on compact convex bodies
$\cA \subset \R^{n_0}$, $\cB \subset \R^{n_1}$. The effective landscapes are
$Q_\nu(a) = \int u(a,b)\,d\nu(b)$ and $W_\mu(b) = -\int u(a,b)\,d\mu(a)$; both
players maximise. $\Lambda_s := \Lambda * \cN(0,\diag(s^2))$. A mixed Nash
equilibrium $(\mu^\star,\nu^\star)$ is fixed and assumed purely atomic with
finite support, $\mu^\star = \sum_{j\le M_0} w^\star_j \delta_{p_j}$,
$\nu^\star = \sum_{l\le M_1} v^\star_l \delta_{r_l}$, all weights positive.

\paragraph{Policy and update.} Player $i$ carries $K_i$ diagonal Gaussian
components with means $m^{(i)}_k$, log-scales $\rho^{(i)}_k = \log\sigma^{(i)}_k$
clipped to $\Sigma_i = [s^{\min}_i, s^{\max}_i]$ coordinatewise, and categorical
logits $\ell^{(i)}$. Writing $\Sigma_k := \diag\big((\sigma^{(i)}_k)^2\big)$ and
using \GC{(2.13)}
\begin{equation}\label{eq:Qreg}
  Q^{\mathrm{reg}}_{\nu,s,\bar a}(a) \;=\; (Q_\nu)_s(a)
  \;-\; \frac{\tau^{\mathrm g}_0}{2}\|a - \bar a\|^2_{\bar\sigma^{-2}},
\end{equation}
\GC{Lemma~1}(ii) puts the Gaussian-head update of player $0$ in the form we use
throughout:
\begin{equation}\label{eq:meanupd}
  m^{(0),+}_k \;=\; m^{(0)}_k \;+\; \eta_0\, w^{(0)}_k\, \Sigma_k\,
  \nabla_a Q^{\mathrm{reg}}_{\nu,\sigma_k,\bar a_k}\big(m^{(0)}_k\big),
\end{equation}
\begin{equation}\label{eq:scaleupd}
  \rho^{(0),+}_{k,d} \;=\; \Big[\rho^{(0)}_{k,d}
  \;+\; \tfrac{\eta_0}{2}\, w^{(0)}_k\, \sigma_{k,d}^2\,
  \partial^2_d (Q_\nu)_{\sigma_k}\big(m^{(0)}_k\big)
  \;+\; \tfrac{\eta_0}{2}\tau^{\mathrm g}_0 w^{(0)}_k
        \big(1 - \sigma_{k,d}^2/\bar\sigma_{k,d}^2\big)\Big]_{\Sigma_0},
\end{equation}
$[\cdot]_{\Sigma_0}$ denoting the clip to the scale range; the identity
$\partial_{(s_d^2/2)}\Lambda_s = \partial_d^2 \Lambda_s$ (the heat equation, as in
\GC{Lemma~3}) has been used. The categorical head runs \GC{(2.6)}, which by
\GC{Lemma~1}(i) is verbatim tabular MMD on the component matrix
$\bar A_{kl}(m,\sigma)$. Player $1$'s equations are the same with $W_\mu$ in place
of $Q_\nu$.

\paragraph{Class.} $u \in \mathfrak{G}(M,\Sigma,\cN,\rho)$ means \GC{Definition~1}:
conditions (C) uniform attracting cell structure with branch maps $\xi_j$,
(D) branch regularity and small gain, (M) modal match, (B) atom regularity,
(I) nondegeneracy of the atom matrix game $A_{jl} = u(p_j,r_l)$, (K) capacity and
covering, (S) algorithmic caps, (W) weight non-degeneracy. We use all of them.

\paragraph{Standing regularity.} By \GC{(2.4)}, for every $k \ge 1$ there is
$C_k$ with $\|\nabla^k (Q_\nu)_s\|_\infty \le C_k \|u\|_\infty / (s^{\min}_0)^k$
uniformly in $\nu \in \cN_1$ and $s \in \Sigma_0^{n_0}$. All constants below are
therefore finite; the floor $s^{\min}_i > 0$ is used for exactly this.

\section{The constants}\label{sec:constants}

Everything in the proof is expressed in the following constants. They are stated
for player $0$; player $1$'s are the same objects for $W_\mu$, with subscript $1$.
All suprema are over $j \le M_0$, $a \in \cA_j$, $\nu \in \cN_1$,
$s \in \Sigma_0^{n_0}$, $\bar a \in \cA_j$, unless said otherwise, and all are
finite by the standing regularity.

\begin{center}
\begin{tabular}{@{}l p{0.72\textwidth}@{}}
\hline
symbol & meaning\\
\hline
$\alpha_0 > 0$ & linear star-concavity modulus, hypothesis~(L) below:
  $\rho(r) \ge \alpha_0 r$ on $(0,\, \mathrm{diam}\,\cA]$.\\
$L_0$ & $\sup \|\nabla^2_a Q^{\mathrm{reg}}_{\nu,s,\bar a}\|_{\mathrm{op}}$, a
  Lipschitz constant for $\nabla_a Q^{\mathrm{reg}}$. Finite; $\le C_2\|u\|_\infty
  /(s^{\min}_0)^2 + \tau^{\mathrm g}_0/(s^{\min}_0)^2$.\\
$L^{(3)}_0$ & $\sup \|\nabla_a \Delta_a (Q_\nu)_s\|$.\\
$\kappa_0 \ge 0$ & uniform branch curvature:
  $\partial_d^2 (Q_\nu)_s(\xi_j) \le -\kappa_0$ for every coordinate $d$. See
  Lemma~\ref{lem:curv}: $\kappa_0 \ge \alpha_0 - \tau^{\mathrm g}_0/(s^{\min}_0)^2$,
  and $\kappa_0 = \alpha_0$ when $\tau^{\mathrm g}_0 = 0$.\\
$w_{\min,0}$ & the weight floor of (W); $\bar w_0 := \max_k w^{(0)}_k \le 1$.\\
$\chi_0$ & anisotropy ratio $\sup_{t,k} \bar\sigma_k/\underline\sigma_k$ where
  $\underline\sigma_k = \min_d \sigma_{k,d}$, $\bar\sigma_k = \max_d\sigma_{k,d}$.
  Equals $1$ for isotropic scales; bounded by $s^{\max}_0/s^{\min}_0$ always.\\
$g^m_0$ & $\Lip$ of $\xi_j$ in the opponent's component means (max over the
  opponent's components, $\ell_\infty$ on the index).\\
$g^w_0$ & $\Lip$ of $\xi_j$ in the opponent's categorical weights ($\ell_1$).\\
$g^\varsigma_0$ & $\Lip$ of $\xi_j$ in the opponent's scale variables
  $(\sigma^{(1)}_l)^2/2$.\\
$g^s_0$ & $\Lip$ of $\xi_j$ in its own scale variable $s^2/2$. By the implicit
  function theorem at the branch, $g^s_0 \le L^{(3)}_0/\alpha_0$
  (Lemma~\ref{lem:branchdrift}).\\
$g^{\bar a}_0$ & $\Lip$ of $\xi_j$ in the magnet centre $\bar a$; $= 0$ when
  $\tau^{\mathrm g}_0 = 0$, and $\le \tau^{\mathrm g}_0/
  (\bar\sigma^2\alpha_0)$ in general.\\
$R_0$ & invariance radius: $\inf_{j,\nu,s,\bar a}
  \mathrm{dist}\big(\xi_j(\nu,s,\bar a),\, \partial\cA_j\big) > 0$.\\
$\Lambda^{A}_0$ & $\sup_{k,l}\|\partial \bar A_{kl}/\partial m^{(0)}_k\|
  = \sup \|\nabla (Q_{\nu_l})_{\sigma_k}(m_k)\|$, the sensitivity of the
  component matrix to player $0$'s means.\\
$\tau_i$ & shorthand for $\tau^{\mathrm{cat}}_i > 0$.\\
$c^-_i,\, c^{\mathrm{qre}}_i,\, c^v_i$ & categorical block constants of
  Lemma~\ref{lem:weight}: contraction rate to the QRE, Lipschitz constant of the
  QRE map in the payoff matrix ($\le C/\tau_i$), and per-step weight displacement
  per unit weight error.\\
$\beta_i$ & the magnet brake factor $\big(1 - e^{-\eta_i\tau^{\mathrm g}_iT}\big)/
  (\eta_i\tau^{\mathrm g}_iT) \in (0,1]$ of \GC{Lemma~6}.\\
\hline
\end{tabular}
\end{center}

\section{The additional hypotheses}\label{sec:hyp}

Definition~\GC{1} is qualitative in three places: $\rho$ is only required
positive, (D) only requires the round-trip gain of the \emph{mean} blocks to be a
contraction, and (S) only requires the step sizes ``below the threshold set by
the Lipschitz constant''. The following name the missing quantities. (L) is
substantive; (D$^+$), (G) and (S$^+$) are the quantitative forms of hypotheses
already present.

\begin{conditions}
\item[\textbf{(L) Linear modulus.}] $\displaystyle
  \alpha_0 \;:=\; \inf_{0 < r \le \mathrm{diam}\,\cA} \frac{\rho(r)}{r} \;>\; 0$,
  and likewise $\alpha_1 > 0$ for player $1$.

\item[\textbf{(D$^+$) Quantified branch regularity.}] The constants
  $g^m_i, g^w_i, g^\varsigma_i, g^s_i, g^{\bar a}_i$ of
  Section~\ref{sec:constants} are finite. (This is (D), with the Lipschitz
  dependence on $\nu$ resolved into its three arguments --- opponent means,
  opponent weights, opponent scales --- rather than left implicit in a metric for
  weak convergence.)

\item[\textbf{(G) Small gain.}] The $4 \times 4$ nonnegative matrix
  $\hat\Gamma$ of \eqref{eq:Gamma} below, indexed by the blocks
  $(r_0, r_1, \delta_0, \delta_1)$, has spectral radius
  $\varrho(\hat\Gamma) < 1$. It suffices that the magnet self-gains obey
  \begin{equation}\label{eq:selfgain}
    \frac{g^{\bar a}_i\, \bar w_i\, L_i\, \chi_i^3}{w_{\min,i}\, \alpha_i}
    \;<\; 1 \qquad (i = 0,1),
  \end{equation}
  and every simple cycle of $\hat\Gamma$ has product $< 1$; the two cycles that
  bind are
  \begin{equation}\label{eq:cyclemean}
    \frac{g^m_0\,g^m_1\,\bar w_0\,\bar w_1\,L_0\,L_1\,\chi_0^3\chi_1^3}
         {w_{\min,0}\,w_{\min,1}\,\alpha_0\,\alpha_1} \;<\; 1
    \tag{mean\,$\leftrightarrow$\,mean}
  \end{equation}
  and, for $\{i,i'\} = \{0,1\}$,
  \begin{equation}\label{eq:cycleweight}
    \frac{g^w_i\, c^v_{i'}\, c^{\mathrm{qre}}_{i'}\, \Lambda^A_i\,
          \bar w_i\, L_i\, \chi_i^3}
         {c^-_{i'}\, w_{\min,i}\, \alpha_i} \;<\; 1 .
    \tag{mean\,$\leftrightarrow$\,weight}
  \end{equation}

\item[\textbf{(G$_\sigma$) Scale-drift dominance.}] With
  \begin{equation}\label{eq:barr}
    \bar r_i \;:=\; \frac{2}{\epsilon}\cdot
    \frac{g^s_i\, \bar w_i\, L_i\, (s^{\max}_i)^2\, \chi_i^2
        \;+\; g^\varsigma_i\, \bar w_{i'}\, L_{i'}\, (s^{\max}_{i'})^2\,
              \chi_{i'}^2 \,\pi_{i'}/\pi_i}
         {2\, w_{\min,i}\, \alpha_i}
  \end{equation}
  ($\epsilon$ and $\pi$ from Lemma~\ref{lem:sg}, $\{i,i'\} = \{0,1\}$), one has
  \begin{equation}
    \bar r_i \;<\; \min\Big\{\, \tfrac{1}{2}R_i,\;\; \delta^\kappa_i \,\Big\},
    \qquad
    \delta^\kappa_i \;:=\; \frac{\kappa_i}{2 L^{(3)}_i},
  \end{equation}
  and the initial means satisfy $\|m^{(i)}_k(0) - \xi_{j(k)}(0)\| < R_i - \bar r_i$
  (the quantitative form of the covering clause of (K)).

\item[\textbf{(S$^+$) Explicit step sizes.}] For $i = 0, 1$,
  \begin{equation}\label{eq:step}
    \eta_i \;<\; \frac{w_{\min,i}\,\alpha_i}
                      {\bar w_i^2\, (s^{\max}_i)^2\, L_i^2\, \chi_i^2},
  \end{equation}
  which makes the second-order term of Lemma~\ref{lem:contract} at most half the
  first-order one.

\item[\textbf{(T) Refresh separation.}] The magnet period $T$ satisfies
  $\eta_i\tau_i T \ge 1$: the categorical block contracts by a definite factor
  between refreshes.
\end{conditions}

\begin{remark}[What (L) costs, and what it buys]\label{rem:L}
(L) is the only hypothesis here that can exclude a game, and it excludes very
little. On a bounded cell, $\alpha_0 > 0$ is equivalent to
$\liminf_{r \to 0}\rho(r)/r > 0$ together with $\rho > 0$ on compacts, i.e.\ to
\emph{nondegeneracy of the branch maximum}, uniformly in $(\nu,s,\bar a)$. It
does not require $\alpha_0$ to be large: for the codebase's Gaussian-bump
landscapes $\rho(r) \asymp (h/w^2)\,r\,e^{-r^2/2w^2}$, so
$\alpha_0 \asymp (h/w^2)e^{-\mathrm{diam}^2/2w^2}$ --- exponentially small, exactly
the ``dead zone is a rate, not a failure'' reading of $s_{\mathrm{reach}}$ in
\GC{\S3.2}. Every conclusion below is uniform in $\alpha_0 > 0$; only the
\emph{rate} degrades as $\alpha_0 \to 0$.

What (L) buys is two things at once: it turns the Lyapunov estimate into a
geometric contraction (Lemma~\ref{lem:contract}), and it supplies the uniform
branch curvature $\kappa_0$ that the scale block needs
(Lemma~\ref{lem:curv}) --- which would otherwise have to be assumed separately.
\end{remark}

\begin{remark}[Why (G) is a legitimate closure and not a restatement]
\label{rem:G}
The point of the constants in \eqref{eq:cyclemean}--\eqref{eq:cycleweight} is
that \emph{no step size and no scale appears in them}. The block contraction
rates are $\eta_i \sigma^2 w \alpha$ and the drifts injected between blocks are
$\eta_{i'}\sigma'^2 w' L$; both $\eta$ and $\sigma^2$ therefore appear once in a
numerator and once in a denominator around any cycle and cancel identically
(Lemma~\ref{lem:scalefree}). This is why (G) can be checked once, offline, from
game data --- and it is the reason the objection of \GC{\S5.6}, item~5 (``the
scale collapse throttles transport by four orders of magnitude'') does not
threaten the argument: it throttles the drift by the same four orders.
\end{remark}

\begin{remark}[$\tau^{\mathrm{cat}}$ is free]\label{rem:tau}
\eqref{eq:cycleweight} is $\propto 1/\tau_{i'}$, since
$c^{\mathrm{qre}}_{i'} \le C/\tau_{i'}$: the mean\,$\leftrightarrow$\,weight cycle
is made a contraction by taking the categorical magnet temperature
\emph{large}. \GC{\S5.6} item~1 records this as an unresolved tension, on the
grounds that large $\tau^{\mathrm{cat}}$ should cost bias. It does not.
Lemma~\ref{lem:weight}(iii) below is the reason: the magnet-refresh outer loop is
a Bregman proximal-point iteration on the matrix game, whose fixed point is a
Nash equilibrium of $\bar A$ for \emph{every} $\tau^{\mathrm{cat}} > 0$. So
$\tau^{\mathrm{cat}}$ enters the convergence \emph{rate} of the outer loop and
not the limit, and may be chosen as large as \eqref{eq:cycleweight} demands. The
ordering of limits that \GC{\S5.6} could not settle is: pick $\tau^{\mathrm{cat}}$
from \eqref{eq:cycleweight}, then $T$ from (T), then $\eta$ from
\eqref{eq:step} --- in that order, with no circularity.
\end{remark}

\section{Preparatory lemmas}\label{sec:lemmas}

\subsection{Regularity at the branch point}

\begin{lemma}[Branch points are critical, and the gradient is dominated]
\label{lem:dom}
Assume (C). For every $j$, $\nu \in \cN_1$, $s \in \Sigma_0^{n_0}$ and
$\bar a \in \cA_j$:
\begin{enumerate}
\item[(i)] $\nabla_a Q^{\mathrm{reg}}_{\nu,s,\bar a}(\xi_j) = 0$, and $\xi_j$ is
  the unique critical point of $Q^{\mathrm{reg}}_{\nu,s,\bar a}$ in $\cA_j$ and
  its unique maximiser over $\overline{\cA_j}$;
\item[(ii)] $\big\|\nabla_a Q^{\mathrm{reg}}_{\nu,s,\bar a}(a)\big\|
  \le L_0 \|a - \xi_j\|$ for every $a \in \cA_j$.
\end{enumerate}
\end{lemma}

\begin{proof}
(i) $\cA_j$ is open and $\xi_j \in \cA_j$, so for small $\epsilon > 0$ and any
unit $v$ the point $a = \xi_j + \epsilon v$ lies in $\cA_j$, and \GC{(3.1)}
reads $\ip{\nabla Q^{\mathrm{reg}}(\xi_j + \epsilon v)}{-\epsilon v} \ge
\rho(\epsilon)\epsilon > 0$, i.e.\
$\ip{\nabla Q^{\mathrm{reg}}(\xi_j + \epsilon v)}{v} \le -\rho(\epsilon) < 0$.
Letting $\epsilon \downarrow 0$ and using continuity of $\nabla Q^{\mathrm{reg}}$
gives $\ip{\nabla Q^{\mathrm{reg}}(\xi_j)}{v} \le 0$ for every $v$, hence
$\nabla Q^{\mathrm{reg}}(\xi_j) = 0$. If $a \ne \xi_j$ were another critical
point, \GC{(3.1)} at $a$ would give $0 \ge \rho(\|\xi_j - a\|)\|\xi_j - a\| > 0$.
For the maximiser claim, integrate \GC{(3.1)} along the segment
$a_t = \xi_j + t(a - \xi_j)$, $t \in (0,1]$, which lies in $\cA_j$ whenever
$\cA_j$ is star-shaped about $\xi_j$ --- and it is, because \GC{(3.1)} makes the
radial derivative of $Q^{\mathrm{reg}}$ toward $\xi_j$ strictly positive on all of
$\cA_j$, so the flow of $\nabla Q^{\mathrm{reg}}$ from any $a \in \cA_j$ reaches
$\xi_j$ without leaving $\cA_j$; along that flow $Q^{\mathrm{reg}}$ is strictly
increasing, so $Q^{\mathrm{reg}}(a) < Q^{\mathrm{reg}}(\xi_j)$. Continuity extends
the inequality to $\overline{\cA_j}$ with $\le$.

(ii) $\nabla Q^{\mathrm{reg}}$ is $L_0$-Lipschitz on $\R^{n_0}$ by the definition
of $L_0$ (the smoothed part by \GC{(2.4)}, the quadratic part exactly), and
vanishes at $\xi_j$ by (i).
\end{proof}

\begin{lemma}[(L) gives uniform curvature at the branch]\label{lem:curv}
Assume (C) and (L). Then for all $(\nu,s,\bar a)$,
\[
  \nabla^2_a Q^{\mathrm{reg}}_{\nu,s,\bar a}(\xi_j) \;\preceq\; -\alpha_0 I,
  \qquad\text{hence}\qquad
  \partial^2_d (Q_\nu)_s(\xi_j) \;\le\; -\alpha_0
  \;+\; \frac{\tau^{\mathrm g}_0}{\bar\sigma_{d}^2}
  \;=:\; -\kappa_0 \quad \text{for every } d .
\]
In particular $\kappa_0 = \alpha_0 > 0$ when $\tau^{\mathrm g}_0 = 0$, and
$\kappa_0 > 0$ whenever the landscape rather than the magnet supplies the
curvature, $\alpha_0 > \tau^{\mathrm g}_0/\bar\sigma_d^2$.
\end{lemma}

\begin{proof}
Put $f = Q^{\mathrm{reg}}_{\nu,s,\bar a}$, $\xi = \xi_j$, and let $v$ be a unit
vector, $a = \xi + \epsilon v$. By Lemma~\ref{lem:dom}(i) and a Taylor expansion,
$\nabla f(a) = \epsilon \nabla^2 f(\xi) v + o(\epsilon)$, so \GC{(3.1)} with (L)
gives
\[
  -\epsilon^2 \ip{\nabla^2 f(\xi) v}{v} + o(\epsilon^2)
  \;=\; \ip{\nabla f(a)}{\xi - a}
  \;\ge\; \rho(\epsilon)\epsilon \;\ge\; \alpha_0 \epsilon^2 .
\]
Dividing by $\epsilon^2$ and letting $\epsilon \downarrow 0$ yields
$\ip{\nabla^2 f(\xi)v}{v} \le -\alpha_0$. The second claim is
$\nabla^2 f = \nabla^2 (Q_\nu)_s - \tau^{\mathrm g}_0 \diag(\bar\sigma^{-2})$
applied to $v = e_d$.
\end{proof}

\begin{lemma}[Branch drift in the scale variable is $O(1)$, not $O(s^{-2})$]
\label{lem:branchdrift}
Assume (C), (L), and that $(\nu,s,\bar a) \mapsto \xi_j$ is differentiable.
Writing $t_d = s_d^2/2$,
\[
  \Big\|\frac{\partial \xi_j}{\partial t_d}\Big\|
  \;\le\; \frac{L^{(3)}_0}{\alpha_0} \;=:\; g^s_0 .
\]
\end{lemma}

\begin{proof}
Differentiate $\nabla_a Q^{\mathrm{reg}}(\xi_j(t)) = 0$ in $t_d$. Since
$\partial_{t_d}(Q_\nu)_s = \partial_d^2 (Q_\nu)_s$ (heat equation) and the
magnet term does not depend on $s$,
\[
  \nabla^2_a Q^{\mathrm{reg}}(\xi_j)\, \frac{\partial\xi_j}{\partial t_d}
  \;+\; \nabla_a\, \partial_d^2 (Q_\nu)_s(\xi_j) \;=\; 0 ,
\]
and $\nabla^2_a Q^{\mathrm{reg}}(\xi_j) \preceq -\alpha_0 I$ is invertible with
$\|(\nabla^2)^{-1}\| \le 1/\alpha_0$ by Lemma~\ref{lem:curv}.
\end{proof}

\begin{remark}
Lemma~\ref{lem:branchdrift} is small but load-bearing. Measured in the
\emph{log}-scale variable $\rho = \log s$ that the algorithm actually updates,
the branch drift is $\|\partial\xi_j/\partial\rho_d\| = s_d^2\,
\|\partial \xi_j/\partial t_d\| \le s_d^2 g^s_0$: it carries a factor $s^2$. That
factor is the same one carried by the Fisher preconditioner in
\eqref{eq:meanupd}, and their cancellation is what makes
Lemma~\ref{lem:scalefree} work.
\end{remark}

\subsection{The mean block}

\begin{lemma}[One-step contraction toward a frozen branch point]
\label{lem:contract}
Assume (C), (L), (W) and \eqref{eq:step}. Fix a component $k$ of player $0$ in
cell $\cA_j$, let $\xi = \xi_j(\nu,\sigma_k,\bar a_k)$ be the branch point for the
\emph{current} $(\nu,\sigma_k,\bar a_k)$, and set
$\tilde r_k := \|m_k - \xi\|_{\Sigma_k^{-1}}$. Then $m_k^+$ of
\eqref{eq:meanupd} satisfies
\begin{equation}\label{eq:contract}
  \|m^+_k - \xi\|_{\Sigma_k^{-1}}
  \;\le\; \Big(1 - \tfrac{1}{2}\eta_0\, w_{\min,0}\, \alpha_0\,
  \underline\sigma_k^2\Big)\, \tilde r_k .
\end{equation}
Moreover the Euclidean displacement obeys
\begin{equation}\label{eq:vel}
  \|m^+_k - m_k\| \;\le\; \eta_0\, \bar w_0\, \bar\sigma_k^{3}\, L_0\, \tilde r_k .
\end{equation}
\end{lemma}

\begin{proof}
Write $d = m_k - \xi$, $G = \nabla_a Q^{\mathrm{reg}}(m_k)$, $w = w^{(0)}_k \in
[w_{\min,0}, \bar w_0]$, and $V = \tfrac12\|d\|^2_{\Sigma_k^{-1}}$. Then
\[
  \tfrac12\|m^+_k - \xi\|^2_{\Sigma_k^{-1}}
  = \tfrac12\|d + \eta_0 w \Sigma_k G\|^2_{\Sigma_k^{-1}}
  = V + \eta_0 w \ip{G}{d} + \tfrac{\eta_0^2 w^2}{2}\|G\|^2_{\Sigma_k} .
\]
By (C) and (L), $\ip{G}{d} = -\ip{G}{\xi - m_k} \le -\rho(\|d\|)\|d\| \le
-\alpha_0\|d\|^2$. By Lemma~\ref{lem:dom}(ii), $\|G\|^2_{\Sigma_k} \le
\bar\sigma_k^2 L_0^2 \|d\|^2$. Also $\|d\|^2 \ge \underline\sigma_k^2
\|d\|^2_{\Sigma_k^{-1}} = 2\underline\sigma_k^2 V$ and $\|d\|^2 \le
\bar\sigma_k^2\|d\|^2_{\Sigma_k^{-1}} = 2\bar\sigma_k^2 V$. Hence
\[
  \tfrac12\|m^+_k - \xi\|^2_{\Sigma_k^{-1}}
  \;\le\; V\Big(1 - 2\eta_0 w_{\min,0}\alpha_0\underline\sigma_k^2
  + \eta_0^2 \bar w_0^2 \bar\sigma_k^4 L_0^2\Big).
\]
By \eqref{eq:step}, $\eta_0 \bar w_0^2\bar\sigma_k^4L_0^2 <
w_{\min,0}\alpha_0 \underline\sigma_k^2$ (using
$\bar\sigma_k^4 \le (s^{\max}_0)^2\chi_0^2\underline\sigma_k^2$), so the bracket is
at most $1 - \eta_0 w_{\min,0}\alpha_0\underline\sigma_k^2$; taking square roots
and using $\sqrt{1-x} \le 1 - x/2$ gives \eqref{eq:contract}. For
\eqref{eq:vel}, $\|m^+_k - m_k\| = \eta_0 w\|\Sigma_k G\| \le \eta_0\bar w_0
\bar\sigma_k^2 L_0\|d\| \le \eta_0\bar w_0\bar\sigma_k^3 L_0 \tilde r_k$.
\end{proof}

\begin{lemma}[Cell invariance]\label{lem:inv}
Under the hypotheses of Lemma~\ref{lem:contract}, if at every step the
\emph{moving} tracking error satisfies $\|m_k(t) - \xi_j(t)\| \le R_0$, then
$m_k(t) \in \cA_j$ for all $t$: the cell assignment $k \mapsto j(k)$ fixed at
initialisation never changes.
\end{lemma}

\begin{proof}
$\|m_k - \xi_j\| \le R_0 \le \mathrm{dist}(\xi_j, \partial\cA_j)$ places $m_k$ in
the ball $B(\xi_j, R_0) \subseteq \cA_j$ by the definition of $R_0$.
\end{proof}

\begin{lemma}[Mean block, moving target]\label{lem:mean}
Under (C), (L), (D$^+$), (W), \eqref{eq:step}, and while $m_k(t) \in \cA_j$, the
quantity $\tilde r_0(t) := \max_k \|m^{(0)}_k(t) - \xi_{j(k)}(t)\|_{\Sigma_k^{-1}}$
obeys
\begin{align}
  \tilde r_0(t+1) \;\le\;& \big(1 - \eta_0 a_0\big)\, \tilde r_0(t)
  \;+\; \frac{g^m_0}{\underline\sigma_0}\, \eta_1 \bar w_1\bar\sigma_1^3 L_1\,
        \tilde r_1(t)
  \;+\; \frac{g^w_0}{\underline\sigma_0}\, \eta_1 \tau_1 c^v_1\, \delta_1(t)
  \nonumber\\
  &\;+\; \frac{1}{\underline\sigma_0}\Big[
        g^s_0\,\big|\Delta (\sigma_0^2/2)\big|
      + g^\varsigma_0\,\big|\Delta (\sigma_1^2/2)\big| \Big],
  \label{eq:meaniss}
\end{align}
where $a_0 := \tfrac12 w_{\min,0}\alpha_0\underline\sigma_0^2
- g^{\bar a}_0\bar w_0\bar\sigma_0^3 L_0 / \underline\sigma_0$, $\underline\sigma_0
= \min_k \underline\sigma_k$, $\bar\sigma_0 = \max_k\bar\sigma_k$, and $\delta_1$
is the categorical error of Lemma~\ref{lem:weight}.
\end{lemma}

\begin{proof}
Insert the frozen-target step of Lemma~\ref{lem:contract} and then move the
target:
\[
  \|m^+_k - \xi^+_j\|_{\Sigma_k^{-1}}
  \le \|m^+_k - \xi_j\|_{\Sigma_k^{-1}}
    + \|\xi^+_j - \xi_j\|_{\Sigma_k^{-1}}
  \le \big(1 - \tfrac{\eta_0}{2} w_{\min}\alpha_0\underline\sigma_k^2\big)\tilde r_k
    + \frac{\|\xi^+_j - \xi_j\|}{\underline\sigma_k}.
\]
By (D$^+$) the target displacement decomposes over its arguments,
\[
  \|\xi^+_j - \xi_j\| \le g^m_0 \max_l \|m^{(1),+}_l - m^{(1)}_l\|
   + g^w_0 \|v^+ - v\|_1
   + g^\varsigma_0 |\Delta(\sigma_1^2/2)|
   + g^s_0 |\Delta(\sigma_0^2/2)|
   + g^{\bar a}_0 \|\bar a^+_k - \bar a_k\| ,
\]
and the five displacements are bounded by \eqref{eq:vel} for player $1$, by
Lemma~\ref{lem:weight}(iv) for the weights, by the scale update
\eqref{eq:scaleupd} for the two scale terms, and --- for the magnet centre --- by
the observation that $\bar a_k$ is refreshed to $m_k$ every $T$ steps, so the
total displacement of $\bar a_k$ over any window of $T$ steps is at most the total
displacement of $m_k$ over the previous window, i.e.\ at most
$T\eta_0 \bar w_0 \bar\sigma_0^3 L_0 \sup \tilde r_0$; amortised over the window
this is the term absorbed into $a_0$. (For $\tau^{\mathrm g}_0 = 0$,
$g^{\bar a}_0 = 0$ and no amortisation is needed; \eqref{eq:meaniss} then holds
step-by-step rather than window-by-window.)
\end{proof}

\subsection{The scale block}

\begin{lemma}[Scale descent and finite-time arrival at the floor]
\label{lem:scale}
Assume (C), (L), (W), and $\kappa_0 > 0$. If at time $t$ every component of
player $0$ satisfies $\|m_k - \xi_{j(k)}\| \le \delta^\kappa_0 =
\kappa_0/(2L^{(3)}_0)$, then every coordinate of every log-scale strictly
decreases,
\begin{equation}\label{eq:scaledesc}
  \rho^{(0)}_{k,d}(t+1) \;\le\; \Big[\rho^{(0)}_{k,d}(t)
  \;-\; \tfrac{\eta_0}{4}\,\beta_0\, w_{\min,0}\, \kappa_0\, \sigma_{k,d}^2(t)
  \Big]_{\Sigma_0} ,
\end{equation}
and consequently, if the hypothesis persists, $\sigma^{(0)}_{k,d}$ reaches the
floor $s^{\min}_0$ in at most
\[
  T^\sigma_0 \;=\; \frac{4}{\eta_0\beta_0 w_{\min,0}\kappa_0 (s^{\min}_0)^2}
  \, \log\frac{s^{\max}_0}{s^{\min}_0}
\]
steps, after which it is constant and $\Delta(\sigma_0^2/2) \equiv 0$.
\end{lemma}

\begin{proof}
By Lemma~\ref{lem:curv}, $\partial_d^2 (Q_\nu)_{\sigma_k}(\xi_j) \le -\kappa_0$.
By the definition of $L^{(3)}_0$, $|\partial_d^2 (Q_\nu)_{\sigma_k}(m_k) -
\partial_d^2 (Q_\nu)_{\sigma_k}(\xi_j)| \le L^{(3)}_0\|m_k - \xi_j\| \le
\kappa_0/2$, so $\partial_d^2 (Q_\nu)_{\sigma_k}(m_k) \le -\kappa_0/2$. Inserting
this in \eqref{eq:scaleupd} and using $w^{(0)}_k \ge w_{\min,0}$ bounds the
landscape term by $-\tfrac{\eta_0}{4} w_{\min,0}\kappa_0\sigma_{k,d}^2$. The
magnet term of \eqref{eq:scaleupd} is $\tfrac{\eta_0}{2}\tau^{\mathrm g}_0
w_k(1 - \sigma_{k,d}^2/\bar\sigma_{k,d}^2)$, which is a logistic relaxation
toward the snapshot $\bar\sigma_{k,d}$; by \GC{Lemma~6}(ii) its effect over a
refresh cycle is exactly to rescale the descent rate by
$\beta_0 \in (0,1]$, giving \eqref{eq:scaledesc}. The arrival time follows by
summing: $\rho$ falls by at least $\tfrac{\eta_0}{4}\beta_0 w_{\min,0}\kappa_0
(s^{\min}_0)^2$ per step while above the floor, and the total distance to fall is
$\log(s^{\max}_0/s^{\min}_0)$. Once at the floor the clip holds it there, because
the unclipped increment is negative.
\end{proof}

\begin{lemma}[Scale-induced tracking error is $O(\sigma^2)$]\label{lem:scalefree}
The exogenous drift terms in \eqref{eq:meaniss} obey
\[
  \frac{1}{\underline\sigma_0}\Big[
        g^s_0\big|\Delta(\sigma_0^2/2)\big|
      + g^\varsigma_0\big|\Delta(\sigma_1^2/2)\big|\Big]
  \;\le\; \frac{\eta_0}{2}\, \frac{g^s_0 \bar w_0 \bar\sigma_0^4 L_0}
                                   {\underline\sigma_0}
        + \frac{\eta_1}{2}\, \frac{g^\varsigma_0 \bar w_1 \bar\sigma_1^4 L_1}
                                   {\underline\sigma_0} ,
\]
so the residual tracking error they can sustain against the contraction rate
$\eta_0 a_0 \asymp \eta_0 w_{\min,0}\alpha_0\underline\sigma_0^2$ is, in Euclidean
units,
\[
  \bar\sigma_0 \cdot \frac{\text{drift}}{\eta_0 a_0}
  \;\asymp\; \frac{g^s_0 \bar w_0 L_0\, \bar\sigma_0^2 \chi_0^2}
                  {w_{\min,0}\alpha_0}
   \;+\; \frac{\eta_1}{\eta_0}\,
         \frac{g^\varsigma_0 \bar w_1 L_1\, \bar\sigma_1^4 \chi_0^2}
              {w_{\min,0}\alpha_0\,\bar\sigma_0^2} ,
\]
i.e.\ it is $O(\sigma^2)$ and \emph{shrinks} as the scales fall.
\end{lemma}

\begin{proof}
$|\Delta(\sigma_{k,d}^2/2)| = \sigma_{k,d}^2 |\Delta\rho_{k,d}| \le
\sigma_{k,d}^2 \cdot \tfrac{\eta_0}{2} \bar w_0 \sigma_{k,d}^2 L_0$ from
\eqref{eq:scaleupd} with $|\partial_d^2 (Q_\nu)_s| \le L_0$ (and the magnet term
only reduces the increment, by \GC{Lemma~6}). The displayed ratio follows on
dividing by $\eta_0 a_0$ and multiplying by $\bar\sigma_0$ to convert
$\|\cdot\|_{\Sigma^{-1}}$ back to $\|\cdot\|$.
\end{proof}

\begin{remark}[Resolution of \GC{\S5.6}, item~5]\label{rem:gap5}
That item objects that Steps~1 and~2 of the skeleton proof are not sequential,
because the mean contraction rate carries $\sigma^2$ and falls by four orders of
magnitude as the scale collapses. Lemmas~\ref{lem:branchdrift} and
\ref{lem:scalefree} answer it: the branch drift generated by the scale collapse
carries $\sigma^4$ --- one factor $\sigma^2$ from the log-scale to scale change of
variables and one from the Fisher preconditioner in the scale update --- so the
\emph{ratio} of drift to contraction carries $\sigma^2$ and improves. Transport is
not throttled relative to what it is chasing; both slow down together, and the
one that slows faster is the disturbance. The same cancellation, applied to the
opponent's mean velocity, is what makes the cycle gains of (G) scale-free
(Remark~\ref{rem:G}).
\end{remark}

\subsection{The weight block}

Let $\bar A$ be the component matrix \GC{(2.9)}, let
$\mathrm{QRE}_\tau(\bar A, \bar w, \bar v)$ be the unique solution of the
$\tau$-entropy-regularised matrix game with magnet $(\bar w, \bar v)$, and let
$\mathrm{Nash}(\bar A)$ be the set of Nash equilibria of $\bar A$. Write
$\delta_0 := \|w^{(0)} - w^{\mathrm{qre}}\|_1$ for the inner tracking error of
player $0$'s categorical head, and $D_0 := \mathrm{dist}_1(\bar w^{(0)},
\mathrm{Nash}(\bar A))$ for the outer error.

\begin{lemma}[Categorical block]\label{lem:weight}
Assume $\tau_i > 0$ and (T). With $\bar A$ and the magnet held fixed:
\begin{enumerate}
\item[(i)] \emph{(inner contraction)} the iteration \GC{(2.6)} is the Bregman
  proximal step for the $\tau$-regularised bilinear game, hence contracts to
  $\mathrm{QRE}_\tau$ at a linear rate: $\delta_i(t+1) \le
  (1 - \eta_i\tau_i c^-_i)\,\delta_i(t)$ with $c^-_i > 0$ absolute;
\item[(ii)] \emph{(sensitivity)} $\mathrm{QRE}_\tau$ is
  $c^{\mathrm{qre}}_i$-Lipschitz in $\bar A$ with $c^{\mathrm{qre}}_i \le C/\tau_i$,
  $C$ absolute;
\item[(iii)] \emph{(outer loop)} the magnet-refresh map $(\bar w,\bar v) \mapsto
  \mathrm{QRE}_\tau(\bar A, \bar w, \bar v)$ is the resolvent of the monotone
  operator of the bilinear game in the KL geometry; its iteration is Bregman
  proximal point, its fixed points are exactly $\mathrm{Nash}(\bar A)$ for every
  $\tau > 0$, and since the associated variational inequality is polyhedral
  (hence metrically subregular), $D_i(n) \to 0$ at a linear rate in the number
  $n$ of refreshes;
\item[(iv)] \emph{(velocity)} the per-step weight displacement obeys
  $\|w^{(i),+} - w^{(i)}\|_1 \le \eta_i\tau_i c^v_i \delta_i$;
\item[(v)] \emph{(ISS)} for time-varying $\bar A(t)$,
  \[
    \delta_i(t+1) \;\le\; (1 - \eta_i\tau_i c^-_i)\,\delta_i(t)
    \;+\; c^{\mathrm{qre}}_i \big\|\bar A(t+1) - \bar A(t)\big\| .
  \]
\end{enumerate}
\end{lemma}

\begin{proof}
(i)--(ii) are the tabular MMD statements of Sokota et al.\ (2023) applied, via
\GC{Lemma~1}(i), to the finite zero-sum game $\bar A$; the lift of a bilinear
zero-sum game is monotone, and the $\tau$-regularised game is $\tau$-strongly
monotone in the KL geometry, which gives both the linear inner rate and the
$1/\tau$ sensitivity of its unique solution to a perturbation of the payoff.
(iii): a fixed point of the refresh map satisfies $z = \mathrm{QRE}_\tau(\bar
A, z)$, whose first-order condition is exactly the variational inequality of
$\bar A$ --- the regularisation term vanishes at $\bar w = w$ --- so the fixed
points are the Nash equilibria for every $\tau > 0$; convergence of Bregman
proximal point on a monotone operator is standard, and linear convergence follows
from metric subregularity, which holds automatically here because the operator
is polyhedral (Robinson). (iv) The step \GC{(2.6)} moves the logits by
$\eta_i(q - \tau_i\log(w/\bar w))/(1 + \eta_i\tau_i)$, which vanishes at the QRE;
$\ell_1$-Lipschitzness of $\mathrm{softmax}$ on the region $w \ge w_{\min,i}$
supplied by (W) gives the constant $c^v_i$. (v) Combine (i) and (ii) with the
triangle inequality through the moved QRE.
\end{proof}

\subsection{A discrete-time small-gain lemma}

\begin{lemma}[Weighted max-form small gain]\label{lem:sg}
Let $x(t) \in \R^N_{\ge 0}$ satisfy, for all $t$ and $i$,
\[
  x_i(t+1) \;\le\; \big(1 - a_i(t)\big)x_i(t)
  \;+\; \sum_{j \ne i} b_{ij}(t)\, x_j(t) \;+\; d_i(t),
\]
with $a_i(t) \in (0,1)$, $b_{ij}(t), d_i(t) \ge 0$. Suppose the matrix
$\hat\Gamma$ with entries $\hat\Gamma_{ij} := \sup_t b_{ij}(t)/a_i(t)$
($i \ne j$, $\hat\Gamma_{ii} := 0$) has $\varrho(\hat\Gamma) < 1$. Then there are
$\pi \in \R^N_{>0}$ and $\epsilon \in (0,1)$ with $\sum_{j\ne i}
\hat\Gamma_{ij}\pi_j \le (1-\epsilon)\pi_i$ for all $i$, and
$V(t) := \max_i x_i(t)/\pi_i$ satisfies
\[
  V(t+1) \;\le\; \big(1 - \epsilon\, \underline a(t)\big)\,V(t)
  \;+\; \max_i \frac{d_i(t)}{\pi_i},
  \qquad \underline a(t) := \min_i a_i(t) .
\]
Consequently: (a) if $d \equiv 0$ and $\sum_t \underline a(t) = \infty$ then
$V(t) \to 0$; (b) if $\max_i d_i(t)/\pi_i \le \theta\,\epsilon\,\underline a(t)$
for all $t$, then $\limsup_t V(t) \le \theta$, and the ball $\{V \le \theta\}$ is
forward invariant.
\end{lemma}

\begin{proof}
Existence of $\pi,\epsilon$: choose $\epsilon \in (0,1)$ small enough that
$\varrho\big(\hat\Gamma/(1-\epsilon)\big) < 1$, put
$\hat\Gamma' := \hat\Gamma/(1-\epsilon)$ and
$\pi := (I - \hat\Gamma')^{-1}\mathbf{1} = \sum_{n\ge0}(\hat\Gamma')^n\mathbf{1}$,
a convergent Neumann series with $\pi \ge \mathbf{1} > 0$. Then
$\hat\Gamma'\pi = \pi - \mathbf{1} \le \pi$, so
$\hat\Gamma\pi \le (1-\epsilon)\pi$ componentwise. Then
\[
  \frac{x_i(t+1)}{\pi_i}
  \le (1-a_i)\frac{x_i}{\pi_i}
    + \frac{a_i}{\pi_i}\sum_{j\ne i}\hat\Gamma_{ij}\pi_j\frac{x_j}{\pi_j}
    + \frac{d_i}{\pi_i}
  \le (1-a_i)V + (1-\epsilon)a_i V + \frac{d_i}{\pi_i}
  = (1-\epsilon a_i)V + \frac{d_i}{\pi_i},
\]
and $1 - \epsilon a_i \le 1 - \epsilon\underline a$. Claim (a) is
$V(t) \le V(0)\prod_{u<t}(1 - \epsilon\underline a(u)) \le
V(0)e^{-\epsilon\sum \underline a}$. Claim (b): if $V(t) \le \theta$ then
$V(t+1) \le (1-\epsilon\underline a)\theta + \theta\epsilon\underline a = \theta$;
and if $V(t) > \theta$ then $V(t+1) \le V(t) - \epsilon\underline a(V(t)-\theta)$,
which drives $V$ into the ball.
\end{proof}

\begin{remark}
For $N = 2$ with zero diagonal, $\varrho(\hat\Gamma) < 1$ is exactly
$\hat\Gamma_{12}\hat\Gamma_{21} < 1$, so the two-block condition
\eqref{eq:cyclemean} is not merely sufficient. For $N = 4$, ``every simple cycle
has product $< 1$'' is necessary for $\varrho < 1$ and sufficient for the
max-form interconnection; for the sum form used here we take
$\varrho(\hat\Gamma) < 1$ itself as the hypothesis (G), which is a $4\times4$
eigenvalue computation from the constants of Section~\ref{sec:constants}.
\end{remark}

\subsection{The limiting configuration}

\begin{lemma}[Self-consistent configuration: existence, uniqueness, and
$O(s^2)$ location]\label{lem:fp}
Assume (C), (L), (D$^+$), (M), (I) and (G). Fix the scales at the floor
$s^{\min} = (s^{\min}_0, s^{\min}_1)$ and define the \emph{consistency map}
$\cT_{s^{\min}}$ on $\big(\prod_j \cA_j\big) \times \Delta_{M_0} \times
\big(\prod_l \cB_l\big) \times \Delta_{M_1}$ by
\[
  \cT_{s^{\min}}(m, w, n, v) \;=\;
  \Big(\ \big(\xi_j(\nu(n,v), s^{\min}_0, m_j)\big)_j,\ \
  \mathrm{Nash}_0\big(\bar A(m,n,s^{\min})\big),\ \ \cdots \Big),
\]
the omitted entries being player $1$'s two analogues. Then
\begin{enumerate}
\item[(i)] $\cT_{s^{\min}}$ is a contraction in the norm
  $\|(m,w,n,v)\|_\pi := \max\{\|m\|_\infty/\pi_1, \|w\|_1/\pi_3,
  \|n\|_\infty/\pi_2, \|v\|_1/\pi_4\}$, with modulus $1 - \epsilon$, where
  $\pi,\epsilon$ are those of Lemma~\ref{lem:sg}; hence it has a unique fixed
  point $z^\infty_{s^{\min}}$;
\item[(ii)] $\|z^\infty_{s^{\min}} - z^\star\|_\pi = O\big((s^{\min}_0)^2 +
  (s^{\min}_1)^2\big)$, where $z^\star = (p, w^\star, r, v^\star)$ collects the
  equilibrium atoms and weights.
\end{enumerate}
\end{lemma}

\begin{proof}
(i) The Lipschitz constants of the four coordinate maps of $\cT$ are precisely
the numerators $b_{ij}$ of the gain matrix $\hat\Gamma$ of \eqref{eq:Gamma},
divided by the corresponding block ``gain to steady state'' $a_i$: for the mean
coordinates, the steady state of \eqref{eq:meaniss} with the drift terms removed
is $\tilde r_0 = (\sum_j b_{0j}x_j)/a_0$; for the weight coordinates it is
Lemma~\ref{lem:weight}(ii)--(iii). So $\cT$ has $\pi$-weighted Lipschitz constant
$\le \max_i \sum_{j} \hat\Gamma_{ij}\pi_j/\pi_i \le 1-\epsilon$ by
Lemma~\ref{lem:sg}. Banach gives the fixed point; the domain is a product of
compact convex sets mapped into itself by (C) (branch points lie in their cells)
and by the simplex structure.

(ii) By (M), $\xi_j(\nu^\star, s^{\min}_0, \cdot) = p_j + O((s^{\min}_0)^2)$, and
by (I) with $\bar A(p, r, s^{\min}) = A + O((s^{\min})^2)$ --- the $O(s^2)$ being
the smoothing of a $C^2$ landscape at its atoms, (B) --- nondegeneracy is an open
condition, so $\mathrm{Nash}(\bar A)$ is single-valued near $A$ and
$\mathrm{Nash}(\bar A(p,r,s^{\min})) = w^\star + O((s^{\min})^2)$. Hence
$\|\cT_{s^{\min}}(z^\star) - z^\star\|_\pi = O((s^{\min}_0)^2 + (s^{\min}_1)^2)$:
$z^\star$ is an approximate fixed point. The standard contraction estimate
$\|z^\infty - z^\star\| \le \|\cT(z^\star) - z^\star\|/\epsilon$ finishes.
\end{proof}

\begin{lemma}[Exploitability at a self-consistent configuration]\label{lem:exp}
Let $(\mu,\nu)$ be the mixture pair induced by a fixed point of
$\cT_{s^{\min}}$: player $0$'s components sit at the branch points
$\xi_j(\nu, s^{\min}_0)$ with weights a Nash of $\bar A$, and symmetrically for
player $1$. Assume (C), (I) and (B). Then, with $a^\star \in \argmax_\cA Q_\nu$
and $b^\star \in \argmax_\cB W_\mu$,
\[
  \varepsilon(\mu,\nu)
  \;=\; \Big[\max_\cA Q_\nu - \max_\cA (Q_\nu)_{s^{\min}_0}\Big]
      + \Big[\max_\cB W_\mu - \max_\cB (W_\mu)_{s^{\min}_1}\Big]
  \;\le\; \tfrac{1}{2}\big(s^{\min}_0\big)^2 \big|\Delta Q_\nu(a^\star)\big|
      \;+\; \tfrac{1}{2}\big(s^{\min}_1\big)^2 \big|\Delta W_\mu(b^\star)\big|
      \;+\; O(s^4).
\]
Moreover the first bracket is also bounded by $\min_j \big[Q_\nu(p_j) -
(Q_\nu)_{s^{\min}_0}(p_j)\big]$, and at $\nu = \nu^\star$ all $M_0$ of these
quantities coincide.
\end{lemma}

\begin{proof}
Player $0$'s regret is $\max_\cA Q_\nu - \E_\mu Q_\nu$, and
$\E_\mu Q_\nu = \sum_j w_j (Q_\nu)_{s^{\min}_0}(\xi_j)$ by \GC{Lemma~1}. Write
$q_j := (Q_\nu)_{s^{\min}_0}(\xi_j) = (\bar A v)_j$. Since $w$ is a Nash of
$\bar A$ and, by (I) plus Lemma~\ref{lem:fp}(ii), fully mixed, complementary
slackness makes all $q_j$ equal to $\max_j q_j$; hence $\E_\mu Q_\nu = \max_j q_j$.
By Lemma~\ref{lem:dom}(i), $\xi_j$ maximises $(Q_\nu)_{s^{\min}_0}$ over
$\overline{\cA_j}$, and the cells cover $\cA$ up to a null set, so by continuity
$\max_j q_j = \max_\cA (Q_\nu)_{s^{\min}_0}$. This gives the displayed identity.
For the bound, $(Q_\nu)_s \le \max Q_\nu$ pointwise (an average of values of
$Q_\nu$), so the bracket is nonnegative; and evaluating $(Q_\nu)_s$ at any point
$a$ gives $\max_\cA (Q_\nu)_s \ge (Q_\nu)_s(a) = Q_\nu(a) + \tfrac{s^2}{2}\Delta
Q_\nu(a) + O(s^4)$ by (B). Taking $a = a^\star$ gives the stated bound; taking
$a = p_j$ and minimising over $j$ gives the second. At $\nu = \nu^\star$ all $p_j$
are maximisers of $Q_{\nu^\star}$ (\GC{Remark~1}) and, by the same complementary
slackness in the smoothed game, all $(Q_{\nu^\star})_s(p_j)$ agree, so the
deficits agree.
\end{proof}

\begin{remark}[A sharper constant than the one previously stated]
\GC{Theorem~1} gives the constant $\tfrac12\sum_j w^\star_j
|\nabla^2 Q_{\nu^\star}(p_j)|$, i.e.\ the \emph{weighted average} of the per-atom
smoothing deficits. Lemma~\ref{lem:exp} shows the exploitability equals a deficit
at a \emph{single} point --- $\max Q_\nu - \max (Q_\nu)_s$ --- because the
categorical head equalises the cells exactly. The two agree at the equilibrium
(all deficits are then equal), so this is a sharpening off the limit rather than
a correction, but it is the form that generalises: it holds with $\min_j$ in place
of $\sum_j w_j$, and it needs no smoothness away from the maximiser.
\end{remark}

\section{The theorem}\label{sec:thm}

Assemble the four ISS blocks with error vector
$x = (\tilde r_0, \tilde r_1, \delta_0, \delta_1)$. Lemma~\ref{lem:mean} for the
two mean blocks and Lemma~\ref{lem:weight}(v) for the two weight blocks give the
comparison system of Lemma~\ref{lem:sg} with
\[
  a_{r_i} = \eta_i a_i, \qquad a_{\delta_i} = \eta_i\tau_i c^-_i,
\]
and the gain matrix $\hat\Gamma_{ij} = \sup_t b_{ij}(t)/a_i(t)$ equal, in the
block order $(r_0, r_1, \delta_0, \delta_1)$, to
\begin{equation}\label{eq:Gamma}
  \hat\Gamma \;=\;
  \begin{pmatrix}
    0
      & \dfrac{\eta_1 g^m_0 \bar w_1 \bar\sigma_1^3 L_1}
              {\eta_0 a_0 \underline\sigma_0}
      & 0
      & \dfrac{\eta_1 g^w_0 \tau_1 c^v_1}{\eta_0 a_0 \underline\sigma_0}
    \\[10pt]
    \dfrac{\eta_0 g^m_1 \bar w_0 \bar\sigma_0^3 L_0}
          {\eta_1 a_1 \underline\sigma_1}
      & 0
      & \dfrac{\eta_0 g^w_1 \tau_0 c^v_0}{\eta_1 a_1 \underline\sigma_1}
      & 0
    \\[10pt]
    \dfrac{c^{\mathrm{qre}}_0 \Lambda^A_0 \bar w_0 \bar\sigma_0^3 L_0}
          {\tau_0 c^-_0}
      & \dfrac{\eta_1 c^{\mathrm{qre}}_0 \Lambda^A_1 \bar w_1 \bar\sigma_1^3 L_1}
              {\eta_0 \tau_0 c^-_0}
      & 0 & 0
    \\[10pt]
    \dfrac{\eta_0 c^{\mathrm{qre}}_1 \Lambda^A_0 \bar w_0 \bar\sigma_0^3 L_0}
          {\eta_1 \tau_1 c^-_1}
      & \dfrac{c^{\mathrm{qre}}_1 \Lambda^A_1 \bar w_1 \bar\sigma_1^3 L_1}
              {\tau_1 c^-_1}
      & 0 & 0
  \end{pmatrix}.
\end{equation}
The exogenous inputs $d_{r_i}$ are the scale-drift terms of
Lemma~\ref{lem:scalefree}; $d_{\delta_i}$ collects the scale-induced drift of
$\bar A$, of the same order.

\begin{lemma}[The cycle gains are scale-free and step-size-free]
\label{lem:cyclefree}
Every simple cycle of \eqref{eq:Gamma} has a product in which all $\eta_i$ and
all $\sigma$'s cancel, up to the anisotropy ratios $\chi_i$. In particular the
two two-cycles are \eqref{eq:cyclemean} and \eqref{eq:cycleweight}, and the
four-cycle $r_0 \to \delta_1 \to r_1 \to \delta_0 \to r_0$ has product
\[
  \frac{g^w_0\, g^w_1\, c^v_0\, c^v_1\, c^{\mathrm{qre}}_0\, c^{\mathrm{qre}}_1\,
        \Lambda^A_0\,\Lambda^A_1\, \bar w_0\,\bar w_1\, L_0\, L_1\,
        \chi_0^3\chi_1^3}
       {c^-_0\, c^-_1\, w_{\min,0}\,w_{\min,1}\,\alpha_0\,\alpha_1},
\]
which is the product of the two cycles \eqref{eq:cycleweight} and is therefore
$< 1$ whenever they are.
\end{lemma}

\begin{proof}
In $\hat\Gamma_{ij}$ the factor $\eta_j$ appears in the numerator and $\eta_i$ in
the denominator; around a cycle each index appears once as $i$ and once as $j$,
so the $\eta$'s cancel. Likewise the numerator of an edge \emph{out of} a mean
block carries $\bar\sigma_j^3$ and the denominator of an edge \emph{into} it
carries $a_j\underline\sigma_j \asymp w_{\min,j}\alpha_j\underline\sigma_j^3$; the
ratio is $\chi_j^3$. Weight-block edges carry no scale in the denominator and
$\bar\sigma^3$ in the numerator only when their source is a mean block, which
supplies the matching $\underline\sigma^3$ on the next edge. Substituting
$a_i \asymp \tfrac12 w_{\min,i}\alpha_i\underline\sigma_i^2$ (which requires the
magnet self-gain bound \eqref{eq:selfgain}) gives the displayed products.
\end{proof}

\begin{theorem}[Global convergence on $\mathfrak{G}$]\label{thm:main}
Let $u \in \mathfrak{G}(M,\Sigma,\cN,\rho)$ in the sense of \GC{Definition~1},
and assume in addition (L), (D$^+$), (G), (G$_\sigma$), (S$^+$), (T), and
$\kappa_0, \kappa_1 > 0$. Run the two-head update --- \GC{(2.6)} on the
categorical heads, \eqref{eq:meanupd}--\eqref{eq:scaleupd} on the Gaussian heads
--- from any initialisation satisfying the covering clause of (K) in the quantitative form of
(G$_\sigma$). Then:
\begin{enumerate}
\item[(a)] \emph{(cell invariance)} the assignment of components to cells fixed at
  initialisation is preserved for all time;
\item[(b)] \emph{(scales)} there is a finite $T^\sigma := \max_i T^\sigma_i$
  such that for $t \ge t_1 + T^\sigma$ every scale of both players equals its
  floor, where $t_1 = O\big(\log(1/\bar r)/(\epsilon\,\eta\,\underline a)\big)$;
\item[(c)] \emph{(transport and equilibration)} $\tilde r_0(t), \tilde r_1(t),
  \delta_0(t), \delta_1(t) \to 0$ geometrically: for $t \ge t_1 + T^\sigma$,
  \[
    \max\Big\{\tfrac{\tilde r_0(t)}{\pi_1}, \tfrac{\tilde r_1(t)}{\pi_2},
    \tfrac{\delta_0(t)}{\pi_3}, \tfrac{\delta_1(t)}{\pi_4}\Big\}
    \;\le\; C_0 \big(1 - \epsilon\,\underline a\big)^{\,t - t_1 - T^\sigma},
    \qquad
    \underline a = \min_i\Big\{\tfrac{\eta_i w_{\min,i}\alpha_i (s^{\min}_i)^2}{2},
    \ \eta_i\tau_i c^-_i\Big\};
  \]
\item[(d)] \emph{(limit)} consequently $(m,w,n,v)(t) \to z^\infty_{s^{\min}}$, the
  unique self-consistent configuration of Lemma~\ref{lem:fp}; every component mean
  converges to the branch point of its own cell, the component matrix converges to
  $\bar A^\infty = A + O(s^2)$, and the categorical heads converge to its unique
  Nash equilibrium --- which by (I) is $(w^\star, v^\star) + O(s^2)$, and without
  (I) is some Nash of $\bar A^\infty$;
\item[(e)] \emph{(location and bias)} $\|m^\infty_j - p_j\| =
  O\big((s^{\min}_0)^2\big)$ for each $j$ after relabelling, and the induced
  policy pair is an $\varepsilon$-Nash equilibrium of the continuous game with
  \[
    \varepsilon \;\le\; \tfrac12 (s^{\min}_0)^2\, \big|\Delta Q_{\nu^\infty}(a^\star)\big|
    \;+\; \tfrac12 (s^{\min}_1)^2\, \big|\Delta W_{\mu^\infty}(b^\star)\big|
    \;+\; O(s^4)
    \;=\; O\big((s^{\min}_0)^2 + (s^{\min}_1)^2\big).
  \]
\end{enumerate}
\end{theorem}

\begin{proof}
\emph{Step 0 (initial assignment).} By (K) each cell contains at least one
initial mean, and by (G$_\sigma$) each initial mean is within $R_i - \bar r_i$ of
its cell's branch point. Fix the map $k \mapsto j(k)$ accordingly.

\emph{Step 1 (the comparison system).} While all means remain in their cells,
Lemmas~\ref{lem:mean} and \ref{lem:weight}(v) give exactly the hypotheses of
Lemma~\ref{lem:sg} for $x = (\tilde r_0, \tilde r_1, \delta_0, \delta_1)$ with the
gain matrix \eqref{eq:Gamma} and the exogenous inputs
$d = (d_{r_0}, d_{r_1}, d_{\delta_0}, d_{\delta_1})$ of
Lemma~\ref{lem:scalefree}. Hypothesis (G) is $\varrho(\hat\Gamma) < 1$, so
Lemma~\ref{lem:sg} applies and yields $\pi, \epsilon$ and the Lyapunov function
$V = \max_i x_i/\pi_i$. By Lemma~\ref{lem:cyclefree} the entries of $\hat\Gamma$
are bounded uniformly in $t$ --- this is where the scale-freeness is used, and it
is what makes $\sup_t$ in Lemma~\ref{lem:sg} finite even though the rates
$a_i(t)$ themselves vary by orders of magnitude along the run.

\emph{Step 2 (the transient ball, and invariance).} By
Lemma~\ref{lem:scalefree}, $\max_i d_i(t)/\pi_i \le \theta\epsilon\underline a(t)$
with $\theta$ the right-hand side of \eqref{eq:barr}; hypothesis (G$_\sigma$) says
$\theta$ translates to Euclidean tracking errors $\bar r_i < \min\{R_i/2,
\delta^\kappa_i\}$. By Lemma~\ref{lem:sg}(b) the ball $\{V \le \max(V(0),\theta)\}$
is forward invariant and $V$ enters $\{V \le \theta\}$ at some finite $t_1$, with
$t_1 \le \log(V(0)/\theta)/(\epsilon\underline a)$. Throughout, the Euclidean
tracking error stays below $R_i$ --- at $t = 0$ by (G$_\sigma$)'s initialisation
clause and thereafter because $V$ never increases beyond its initial value --- so
Lemma~\ref{lem:inv} gives (a).

\emph{Step 3 (scales, giving (b)).} For $t \ge t_1$ every Euclidean tracking error
is $\le \bar r_i \le \delta^\kappa_i$, so Lemma~\ref{lem:scale} applies at every
step: every log-scale decreases by at least $\tfrac{\eta_i}{4}\beta_i w_{\min,i}
\kappa_i (s^{\min}_i)^2$ per step until it is clipped at the floor, which happens
within $T^\sigma_i$ steps. The clip then holds it there, since the unclipped
increment remains negative (the hypothesis of Lemma~\ref{lem:scale} continues to
hold, by the invariance in Step 2). This proves (b).

\emph{Step 4 (exact convergence, giving (c)).} For $t \ge t_1 + T^\sigma$ all
scales are constant at their floors, so $\Delta(\sigma_i^2/2) \equiv 0$ and the
exogenous inputs vanish: $d(t) \equiv 0$. Lemma~\ref{lem:sg}(a) with the constant
rates $a_{r_i} = \tfrac12 \eta_i w_{\min,i}\alpha_i (s^{\min}_i)^2$ and
$a_{\delta_i} = \eta_i\tau_i c^-_i$ gives the geometric decay stated in (c). This
is the step at which the argument would fail without Lemma~\ref{lem:dom}(ii): if
the mean update's step were merely $O(\eta)$ rather than $O(\eta\,\tilde r)$, the
inputs $b_{ij}x_j$ would not vanish with $x$ and Lemma~\ref{lem:sg} would deliver
only an $O(\eta)$ ball.

\emph{Step 5 (identifying the limit, giving (d)).} $x(t) \to 0$ means the
configuration converges to a point at which every mean is exactly at its own
cell's branch point for the limiting opponent policy, and every categorical head
is exactly at the QRE for the limiting component matrix; by
Lemma~\ref{lem:weight}(iii) and (T), the outer magnet-refresh loop drives that
QRE to a Nash of $\bar A^\infty$ at a linear rate in the number of refreshes, and
the refresh count grows linearly in $t$, so the composite rate is again
geometric. The limit is therefore a fixed point of the consistency map
$\cT_{s^{\min}}$ of Lemma~\ref{lem:fp}; by Lemma~\ref{lem:fp}(i) that fixed point
is unique, so the limit is $z^\infty_{s^{\min}}$. Duplicated components in one
cell converge to the same branch point; the induced matrix then has duplicated
rows and the Nash is non-unique in $w$ but unique as a measure on the branch
points, which is the object the conclusion is stated for.

\emph{Step 6 (location and bias, giving (e)).} Lemma~\ref{lem:fp}(ii) gives
$\|m^\infty_j - p_j\| = O((s^{\min}_0)^2)$ and $\|w^\infty - w^\star\|_1 =
O(s^2)$. Lemma~\ref{lem:exp} applied at $z^\infty_{s^{\min}}$ --- which is a fixed
point of $\cT$, so its hypotheses hold exactly --- gives the exploitability bound.
\end{proof}

\begin{remark}[Which hypothesis does which job]
(C)+(L) do all the work in Steps 2--4 and are the only conditions the negative
result \GC{Proposition~1} forces. (D$^+$)+(G) are used exactly once, to bound
$\varrho(\hat\Gamma)$, i.e.\ to stop the two players' targets from chasing each
other. (G$_\sigma$) is used exactly once, to guarantee that the transient
tracking error is inside the radius where the scale block is monotone. (M)+(I)+(B)
are used only in Step 6, to \emph{name} the limit and give the bias its exponent;
if they fail, (a)--(d) survive verbatim and only (e) is lost. (K) is used only in
Step 0, and (W) only through $w_{\min}$ in Lemma~\ref{lem:contract} --- vacuous
under the corrected preconditioner of \GC{Remark~2}. (S$^+$) and (T) are step-size
bookkeeping.
\end{remark}

\section{Corollaries}\label{sec:cor}

\begin{corollary}[Released floor: transport survives, and $\varepsilon =
\Theta(1/t)$]\label{cor:unfloored}
Suppose the hypotheses of Theorem~\ref{thm:main} hold for every floor
$s^{\min} \in (0, \bar s]$ with constants uniform in $s^{\min}$ --- which they do
whenever $L_i, L^{(3)}_i, \alpha_i, \kappa_i$ are taken as the landscape's own
constants near its atoms rather than the worst-case bounds of \GC{(2.4)} --- and
release the floor. Then:
\begin{enumerate}
\item[(i)] $\sigma_{i}(t) = \big(\sigma_i(0)^{-2} + 2c_it\big)^{-1/2} =
  \Theta(t^{-1/2})$ with $c_i = \tfrac{\eta_i}{2}\beta_i w_i\kappa_i$, so the
  ratio $\sigma_1(t)/\sigma_0(t) \to \sqrt{c_0/c_1}$ is bounded;
\item[(ii)] consequently the entries of $\hat\Gamma$ remain uniformly bounded, so
  (G) still holds and Lemma~\ref{lem:sg} still applies, now with time-varying
  rates $\underline a(t) = \Theta(\eta/t)$;
\item[(iii)] since $\sum_t \underline a(t) = \Theta(\log t) \to \infty$, the
  tracking errors still converge to zero, at the polynomial rate
  $V(t) = O(t^{-\lambda})$ with $\lambda = \epsilon\,\min_i w_{\min,i}\alpha_i
  /(2c_i)$;
\item[(iv)] and $\varepsilon(t) = \Theta(1/t)$, as in \GC{Corollary~2}.
\end{enumerate}
\end{corollary}

\begin{proof}
(i) is \GC{Corollary~2} together with \GC{Lemma~6}; the ratio statement follows
from the closed form. (ii) is Lemma~\ref{lem:cyclefree}: the cycle gains involve
$\sigma$ only through $\chi_i$ and through the ratio $\bar\sigma_1/\underline
\sigma_0$ in individual entries, and (i) bounds the latter. (iii) is
Lemma~\ref{lem:sg}(a) with $\underline a(t) = \eta w_{\min}\alpha\sigma^2(t)/2
= \Theta(1/t)$: $\prod_{u \le t}(1 - \epsilon\underline a(u)) =
\exp(-\epsilon\sum_u \underline a(u)) = \exp(-\lambda\log t + O(1))$. (iv) is Step
6 of Theorem~\ref{thm:main} with $s^{\min}$ replaced by $\sigma(t)$, valid because
by (iii) the tracking error decays polynomially and hence faster than the
$\Theta(1/t)$ bias for $\lambda > 1$; for $\lambda \le 1$ the exploitability is
$\Theta(t^{-\min(1,\lambda)})$.
\end{proof}

\begin{remark}
Corollary~\ref{cor:unfloored} is the second half of the answer to \GC{\S5.6},
item~5. That item worried that unfloored, ``both the contraction and the branch
drift vanish together and the ISS estimate needs their \emph{ratio} controlled''.
It does, and the ratio is controlled: by Lemma~\ref{lem:cyclefree} it is exactly
the scale-free cycle gain, and the only new thing needed is that the two players'
scales fall at comparable rates, which (i) supplies in closed form. The price is
that the convergence of the \emph{means} degrades from geometric to polynomial,
with an exponent $\lambda$ that the user controls through $\eta$ and $\tau^{\rm g}$
(both enter $c_i$).
\end{remark}

\begin{corollary}[Overcapacity]\label{cor:overcap}
Under the hypotheses of Theorem~\ref{thm:main}, taking $K_i \gg M_i$ with a
space-filling initialisation discharges the covering clause of (K) and
(G$_\sigma$) automatically, at no cost to conclusions (a)--(e): spare components
converge to branch points of whichever cell they start in, i.e.\ to genuine atoms,
and the induced measure is unaffected (Step 5).
\end{corollary}

\begin{corollary}[Flat directions: measure convergence only]\label{cor:flat}
Suppose $\kappa_i = 0$ in some directions --- e.g.\ a pure-coupling game, where
$Q_{\nu^\star} \equiv \mathrm{const}$ and (C) holds only via
$\tau^{\mathrm g}_i > 0$ with $\alpha_i = \tau^{\mathrm g}_i/\bar\sigma^2$ and a
single cell (\GC{Remark~5}). Then Steps 1, 2, 4 and 5 of
Theorem~\ref{thm:main} go through unchanged --- the mean and weight blocks are
unaffected --- but Step 3 fails: the scales do not descend from landscape
curvature, and the limit set of Lemma~\ref{lem:fp} is a manifold (any
moment-matched configuration) rather than a point, so $\cT$ is not a contraction
in the parameters. The surviving conclusion is convergence of the \emph{induced
measure} and of the exploitability, with the bias set by whatever scale the
magnet's own relaxation $\sigma \to \bar\sigma$ settles at, and no statement about
the individual means.
\end{corollary}

\begin{proposition}[Closed-form (G) under a finite-rank coupling]
\label{prop:cap}
Let $u(a,b) = D_0(a) - D_1(b) + \ip{\varphi(a)}{C\psi(b)}$ as in \GC{(6.1)}, and
suppose $\nabla^2 D_{i,s} \preceq -\alpha_i I$ near the branches. Then
$g^m_0 \le \|C\|\,\|\nabla\varphi_s\|\,\|\nabla\psi_s\|/\alpha_0$ and
symmetrically, so \eqref{eq:cyclemean} is implied by
\begin{equation}\label{eq:capplus}
  \alpha_0^2\,\alpha_1^2 \;>\;
  \|C\|^2\,\sup\|\nabla\varphi\|^2\,\sup\|\nabla\psi\|^2 \cdot
  \frac{\bar w_0\bar w_1 L_0L_1 \chi_0^3\chi_1^3}{w_{\min,0}w_{\min,1}} .
\end{equation}
This is \GC{Proposition~2}'s cap $\alpha_0\alpha_1 > \|C\|^2\sup\|\nabla\varphi\|
\sup\|\nabla\psi\|$ with the condition numbers $L_i/\alpha_i$ and
$\bar w_i/w_{\min,i}$ made explicit; on \texttt{TWO\_POINT}
($\alpha = 100$, $\|f'\| = 2$, $c = 1$, $L \approx \alpha$, isotropic scales,
$K = M = 2$ so $\bar w/w_{\min} \approx 1$), the left side is $10^8$ and the right
side $\approx 16$, so \eqref{eq:capplus} holds with more than six orders of
margin --- consistent with \GC{Example~1}, which found $g_0g_1 = 4\times10^{-4}$.
\end{proposition}

\begin{proof}
$\xi_j$ solves $\nabla D_{0,s}(\xi) + \nabla\varphi_s(\xi)^\top C y = 0$ with
$y = \E_\nu\psi$; the implicit function theorem with
$\|(\nabla^2 D_{0,s})^{-1}\| \le 1/\alpha_0$ gives
$\|\partial\xi/\partial y\| \le \|\nabla\varphi_s\|\|C\|/\alpha_0$, and
$\|\partial y/\partial m^{(1)}_l\| \le \|\nabla\psi_s\|$ by the chain rule through
the mixture. Substitute in \eqref{eq:cyclemean}.
\end{proof}

\section{What is still not proved}\label{sec:remaining}

Theorem~\ref{thm:main} closes items~1 and~5 of \GC{\S5.6}, and
Corollary~\ref{cor:flat} converts item~4 from ``stated but not proved'' into a
proved but weaker statement. What remains:

\begin{enumerate}
\item \textbf{Exact gradients.} Everything above assumes $U$ and its derivatives
  are available exactly. With stochastic estimates the mean block acquires an
  additive noise term whose variance does not vanish with $\tilde r$, so
  Lemma~\ref{lem:sg} would give an $O(\sqrt{\eta})$ ball rather than zero, and a
  decreasing step size or variance reduction would be needed. This is the single
  largest remaining gap between the theorem and the implementation.
\item \textbf{Verification of (C) and of (L).} Unchanged from \GC{\S5.6} item~3:
  the quantifier over $\nu \in \cN_1$ and $s \in \Sigma_i$ is collapsed by a
  finite-rank coupling (Proposition~\ref{prop:cap}) or by one-dimensional
  scale-space causality, and not otherwise. Note that (L) inherits exactly the
  same checkability: $\alpha_i$ is a mode-ladder computation, not a new one.
\item \textbf{Quantitative $\rho$, i.e.\ $\alpha_i$ in closed form.}
  \GC{\S5.6} item~2. The theorem is now explicitly uniform in $\alpha_i > 0$ and
  every rate is proportional to it, so this has been reduced from ``a missing
  hypothesis'' to ``a missing constant''; computing $\alpha_i$ for
  Gaussian-bump landscapes is elementary and gives
  $\alpha \asymp (h/w^2)e^{-\mathrm{diam}^2/2w^2}$, from which
  Theorem~\ref{thm:main}(c) yields an explicit finite-horizon bound --- the honest
  form of the old $s_{\mathrm{reach}}$ edge.
\item \textbf{The anisotropy ratio $\chi_i$.} It enters the cycle gains as
  $\chi_i^3$, and is bounded a priori only by $s^{\max}_i/s^{\min}_i$. It equals
  $1$ for isotropic scales, and can be removed entirely by imposing (C) in
  preconditioned form (star-concavity of $Q^{\mathrm{reg}}$ measured in the
  $\Sigma_k^{-1}$ metric); whether the resulting condition is still checkable by
  mode counting is not settled.
\item \textbf{The Gaussian magnet's refresh.} Lemma~\ref{lem:mean} amortises the
  magnet-centre jump over a refresh window rather than handling it stepwise, which
  costs a factor $T$ hidden in the constant $C_0$ of Theorem~\ref{thm:main}(c).
  For $\tau^{\mathrm g}_i = 0$ --- the configuration \GC{\S5.4} recommends in
  curved games --- this is vacuous.
\item \textbf{Extensive form.} Unchanged: the payoff is bilinear in sequence form
  but the policy parameterises behavioural strategies per infoset, and the
  cellwise analysis would have to be redone with continuation values as the
  landscape.
\end{enumerate}

\section{Summary}

The class $\mathfrak{G}$ was defined so that the only obstruction to convergence
lives in the mean dynamics. Making that into a theorem needed three things, all
now supplied. The mean update's \emph{step} is proportional to its own tracking
error, because the branch point is a critical point of the regularised landscape
(Lemma~\ref{lem:dom}); this converts input-to-state stability into contraction.
The Fisher preconditioner $\sigma^2$ appears once in every contraction rate and
once in every drift term, so it cancels around every cycle
(Lemma~\ref{lem:cyclefree}); this makes the small-gain condition a property of the
game rather than of the schedule, and voids the ``scale collapse throttles
transport'' objection. And the categorical magnet temperature is free to be taken
large, because the magnet refresh is a proximal-point iteration whose fixed point
is Nash for every $\tau^{\mathrm{cat}} > 0$ (Lemma~\ref{lem:weight}(iii)); this
resolves the ordering of limits that blocked the composition.

The cost is one substantive hypothesis, (L): the star-concavity modulus is linear,
with a constant that may be exponentially small. That constant does not gate
convergence --- it sets the rate. Which is the same statement, in the language of
the class, that \GC{\S3.2} made about the dead zone: it is a rate, not a failure.

\end{document}
